\documentclass[11pt,a4paper]{article}
\usepackage[T1]{fontenc}
\usepackage{lmodern}
\usepackage[margin=24mm,headheight=14pt]{geometry}
\usepackage{amsmath,amssymb,booktabs,longtable,array,microtype}
\usepackage[table]{xcolor}
\usepackage{xurl,hyperref,listings,fancyhdr,enumitem,needspace}
\definecolor{ink}{HTML}{18354A}
\definecolor{accent}{HTML}{007C83}
\definecolor{light}{HTML}{F2F6F8}
\hypersetup{colorlinks=true,linkcolor=ink,urlcolor=accent,
  pdftitle={Further Lean Proof Simplifications for DiophantineConsistency},
  pdfauthor={Prepared for Vladimir Reshetnikov}}
\setlength{\parindent}{0pt}
\setlength{\parskip}{6pt}
\setlength{\emergencystretch}{3em}
\setlist{itemsep=3pt,topsep=4pt,leftmargin=18pt}
\pagestyle{fancy}
\fancyhf{}
\lhead{\small\color{ink}Further Lean proof simplifications}
\rhead{\small 26 October 2026}
\cfoot{\small\thepage}
\lstdefinelanguage{Lean}{keywords={theorem,by,have,exact,def,structure,where,
  namespace,end,import,induction,generalizing,with,fun,forall,let,if,then,else,
  simp,only,apply,refine,change,intro,obtain,rcases},sensitive=true,
  morecomment=[l]{--},morecomment=[s]{/-}{-/},morestring=[b]"}
\lstset{language=Lean,basicstyle=\ttfamily\small,keywordstyle=\color{ink}\bfseries,
  commentstyle=\color{gray},backgroundcolor=\color{light},frame=single,
  rulecolor=\color{light},breaklines=true,columns=fullflexible,keepspaces=true,
  showstringspaces=false,aboveskip=9pt,belowskip=7pt,xleftmargin=5pt,xrightmargin=5pt}
\newcommand{\code}[1]{\texttt{\detokenize{#1}}}
\newcommand{\repo}{https://github.com/VladimirReshetnikov/ProveIt}
\newcommand{\pin}{23f631ea07fff8632eb8cd92286f7621bbe051c0}
\newcommand{\src}[2]{\href{\repo/blob/\pin/Computability/DiophantineConsistency/Lean/DiophantineConsistency/#1.lean}{#2}}
\newcommand{\hsrc}[2]{\href{\repo/blob/\pin/Computability/HilbertTenthProblem/Lean/Diophantine/#1.lean}{#2}}
\newcommand{\fullsrc}[2]{\href{\repo/blob/\pin/#1}{#2}}
\newcommand{\mathsrc}[2]{\href{https://github.com/leanprover-community/mathlib4/blob/81a5d257c8e410db227a6665ed08f64fea08e997/#1}{#2}}
\newcolumntype{P}[1]{>{\raggedright\arraybackslash}p{#1}}
\newcommand{\RawSaved}{223}
\newcommand{\IncrementalSaved}{219}
\newcommand{\RatioSaved}{86}
\newcommand{\RatioRawSaved}{90}
\newcommand{\CombinedModules}{71}
\newcommand{\TreeFiles}{618}
\newcommand{\TreeLines}{80,297}


\begin{document}
\hypersetup{pageanchor=false}
\begin{titlepage}
\vspace*{17mm}
{\color{accent}\large A SOURCE-PINNED, CHECKED PROPOSAL PACKET\par}
\vspace{8mm}
{\color{ink}\Huge\bfseries Further Lean Proof\par Simplifications\par}
\vspace{6mm}
{\LARGE\code{DiophantineConsistency}\par}
\vspace{10mm}
{\large Four additions beyond the existing proposal corpus\par}
\vspace{13mm}
\begin{minipage}{0.94\textwidth}
\large
The strongest change replaces manual proofs of object-formula specifications
with explicit compositional semantic normalization. A second patch shares
one scale-parameterized Pell estimate between ordinary and native consumers.
Two checked architectural options isolate integer recovery and internal
certificate closure.
\end{minipage}
\vfill
\begin{tabular}{@{}ll@{}}
Prepared for & Vladimir Reshetnikov\\[4pt]
Date & 26 October 2026\\[4pt]
Repository & \href{\repo}{VladimirReshetnikov/ProveIt}\\[4pt]
Source commit & \texttt{23f631ea07fff8632eb8cd92286f7621bbe051c0}\\[4pt]
Toolchain & Lean 4.32.0, repository-pinned mathlib\\[4pt]
Deliverables & Report, exact source overlays, patches, examples, evidence
\end{tabular}
\vspace{8mm}
\end{titlepage}
\pagenumbering{roman}
\hypersetup{pageanchor=true}
\tableofcontents
\newpage
\pagenumbering{arabic}

\section{Results and the counterfactual baseline}

\textbf{Recommendation.} Adopt the local specification proofs first, then
review the common Pell estimate and its exact ordinary/native adapters as one
mathematical refactor. Keep integer recovery and observed-set closure as
optional abstractions until their client population justifies the additional
interface.

The two integrated patch families remove \RawSaved{} physical source lines
against the pinned checkout, counting new helpers and imports. The earlier
corpus already proposed simplifying trivial existential aliases inside the
Pell estimates. Charging that prior improvement to the baseline leaves
\IncrementalSaved{} additional lines removed in the explicitly scoped
counterfactual. This adjustment is not a claim to have merged and rebuilt every
alternative design in the earlier packets.

\begin{table}[ht]
\centering\small
\begin{tabular}{@{}P{60mm}P{28mm}P{34mm}P{20mm}@{}}
\toprule
Direction & Concrete scope & Validation level & Net lines\\
\midrule
Compositional formula specifications & 20 proofs; 8 modules & Complete production modules & $-133$\\
One scaled Pell ratio estimate & 3 modules, shared theorem & Complete production modules & $-\RatioSaved$\\
Two-error integer recovery & Core and 3 applications & Checked separate example & Unclaimed\\
Observed certificate-set closure & Core and 6 client theorems & Checked separate examples & Unclaimed\\
\bottomrule
\end{tabular}
\caption{Line reductions include maintained helper costs. The Pell figure is
the incremental reduction after the prior alias cleanup. Optional examples are
not included in the total.}
\end{table}

The requested Lean tree contains \TreeFiles{} files and \TreeLines{} physical
lines at the pin, including comments and blank lines. The deliverable changes
24 existing proofs in 11 production files; its common helper is added to an
existing module. One of the existing files is the
lower-level Hilbert-tenth-problem Pell ratio module: putting the common theorem
there permits both existing consumers to use it without reversing the import
direction. Two optional closure client families also live in imported foundational
modules outside the requested tree; their possible savings are not counted.

The source snapshot is the \href{\repo/commit/\pin}{26 October commit
\texttt{23f631ea07fff863}}. The patch retains public theorem statements and the
existing mathematical objects. It does not change a Diophantine equation,
coefficient, witness count, tile table, encoded formula, or native factor of
two. Source compression here is an engineering result, not a stronger
consistency theorem or a measured compilation-speed improvement.

\section{What was excluded, and why these directions remain}

The audit treats \emph{all} suggestions in the two named locations as already
available: adopted changes, pending patches, uncompiled sketches, deferred
architecture, and supplementary research directions. Reimplementing an old
idea with a new theorem name does not qualify. The supplied
\path{evidence/novelty-audit-initial.txt} records 53 initial excluded families.
The three \path{evidence/late-audit-part*.txt} files extend that review to the
four incoming archives added during this task. Inventories record the archive
and source hashes.

The scope includes the six research packets
\href{\repo/tree/\pin/Computability/DiophantineConsistency/research/lean-proof-simplification}{2026-10-26a through 2026-10-26f},
the nine proof-simplification archives physically present in
\href{\repo/tree/\pin/docs/incoming}{\code{docs/incoming}}, and relevant index
descriptions. The index also links historical mathematics archives elsewhere
in the repository; this report does not claim a semantic review of every
transitively linked historical archive. The incoming ZIPs were read as inert
data; their supplied scripts were not executed.


\textbf{The final refresh changed the recommendations.} The source and all
production candidates were rebased to \texttt{23f631ea07fff863}. The new
\code{additional_lean_simplifications} archive already proposes a stronger
fundamental-unit simplification for the exact client, so the small-square
alternative was withdrawn. Round2's screening entry N08 and New Deltas'
supplementary entry X3 already propose rail/rank reuse, so the finite-cut
family was withdrawn too. Finally, the new-simplifications archive deletes
\code{zeroSumStage_spec} and \code{zeroProductStage_spec} by eliminating
redundant inductions; those two targets were removed from the specification
patch. Tentative status or a different proof does not make a prior idea new.

\begin{longtable}{@{}P{36mm}P{55mm}P{58mm}@{}}
\toprule
New contribution & Nearest already-proposed work & Additional obligation addressed here\\
\midrule\endhead
Formula-specification normalization & Definable-induction wrappers, model-validity bridges, internal recursion preservation, constructor views & Deriving the exact \code{Sat} specification under binders; the earlier wrappers consume it.\\[4pt]
Scaled Pell estimate & Source-register decoding and trivial \code{set}/\code{let} alias cleanup & One analytic error estimate parameterized by a positive scale, with exact ordinary and doubled clients.\\[4pt]
Integer recovery & Half-radix identities and cyclic carry normalization & A unit-interval uniqueness theorem with an explicit sum of two real error budgets.\\[4pt]
Observed certificate closure & Guarded sequential history folds, internal constructor inversion and unique-existence induction & Union and singleton insertion for arbitrary internal sets, observed through their pair/triple encodings.\\
\bottomrule
\end{longtable}

Several tempting changes were therefore rejected: generic sum-of-squares
lemmas, degree automation, circuit certificates, decoder fusion, Pell parity
reuse, functional-graph witness elimination as a broad direction, and another
induction transport wrapper. The current packet uses ordinary validation and
scoped simplification practices, but does not count those practices as new
ideas.

\section{Normalize exact object-formula specifications compositionally}
\label{sec:spec}

\subsection{The repeated proof obligation}

Many formulas are deliberately explicit: quantifier order, de Bruijn slots,
naturalness guards, and represented arithmetic atoms are part of their
meaning. Their specification theorems then prove that interpretation by
\code{Sat} agrees with a semantic predicate. A manual proof often introduces
every binder, splits an equivalence into two directions, constructs the same
existential witnesses, applies an atom's specification using \code{.mp} or
\code{.mpr}, and follows each environment projection separately.

These steps are mostly the recursive equations of satisfaction plus facts
already proved about the atoms. The useful boundary is explicit, local
normalization: unfold the selected outer formula and the logical clauses of
\code{Sat}; rewrite its atomic graphs through their exact specifications;
reduce \code{scons}; finish with only the propositional rules the formula
needs. This does not require a custom tactic or a global simp policy.

\Needspace{9\baselineskip}
For example, the complete replacement body for
\src{ObjectAdditionLaws}{\code{commutativeStage_spec}} is:

\begin{lstlisting}
by
  simp only [commutativeStage, Sat, addForm_spec H, scons]
\end{lstlisting}

The successor-left specification also has determined successor witnesses.
Its checked body makes their elimination explicit:

\begin{lstlisting}
by
  simp only [successorLeftStage, Sat, addForm_spec H,
    fSuccF_spec H, scons, exists_eq_left]
\end{lstlisting}

This is not a suggestion to indiscriminately unfold every object-language
definition. Each proof supplies its own finite list of semantic facts. In
\src{ObjectCodeBound}{\code{boundAtomForm_spec}}, naturalness and tag guards
occur beside two determined successor witnesses. The proof additionally uses
\code{exists_and_left} to expose those equalities and \code{exists_eq_left} to
substitute them. The guards remain in the resulting proposition.

\subsection{A structural induction becomes a semantic fold}

The more substantial example is
\src{ObjectNaturalExpression}{\code{evalForm_spec}}. Its induction is necessary
because the object formula is compiled recursively from a natural expression.
What can disappear is the repeated binder and environment bookkeeping inside
the addition and multiplication cases:

\begin{lstlisting}
by
  induction e generalizing coordinates z ee with
  | constant n => exact fNumF_spec H n z ee
  | var i => rfl
  | add p q hp hq | mul p q hp hq =>
    simp only [evalForm, Sat, EvalRel, hp, hq,
      addForm_spec H, mulForm_spec H, Function.comp_def, scons]
\end{lstlisting}

Generalizing coordinates, the output slot, and the environment is retained:
recursive formulas run beneath additional binders. Removing that
generalization would make the induction hypotheses too weak. The proof still
establishes the semantics of the actual formula compiler, even if a prior
semiring abstraction has already organized its surrounding algebra.

\subsection{Measured scope and preservation}

\begin{center}\small
\begin{tabular}{@{}lrr@{}}
\toprule
Module & Proofs & Lines removed\\
\midrule
\code{ObjectAdditionLaws} & 3 & 20\\
\code{ObjectAdditionCancellation} & 1 & 4\\
\code{ObjectMultiplicationLaws} & 2 & 10\\
\code{ObjectDistributivity} & 2 & 14\\
\code{ObjectNaturalOrder} & 2 & 10\\
\code{ObjectNaturalExpression} & 1 & 17\\
\code{ObjectCodeBound} & 6 & 41\\
\code{ObjectCodeSeal} & 3 & 17\\
\midrule
Total & 20 & 133\\
\bottomrule
\end{tabular}
\end{center}

The eight modules shrink from 1,713 to 1,580 physical lines, removing 5,283
UTF-8 bytes. The selected proof bodies shrink from 185 to 52 lines. No helper,
macro, or global simp declaration is added. A deterministic erasure audit
found every byte outside the selected proof bodies identical, including
theorem statements, formula data, quantifier order, and executable syntax.
All eight complete replacement modules compile together.

The model carrier remains arbitrary. No object-level natural is assumed
standard, and no internal induction is replaced with external induction on
Lean's \code{Nat}. Earlier preservation-stage targets in
\code{ObjectArithmetic}, \code{ObjectMultiplication}, and
\code{ObjectCodeClosure} are deliberately omitted. The separate totality and
uniqueness specifications in \code{ObjectCodeBound} are also omitted because
the earlier unique-existence proposal already targets that layer.
The late induction-elimination proposal also removes the zero-sum and
zero-product stage specifications; those targets are excluded rather than
counted as remaining opportunities.

The principal maintenance risk is a future change in a formula's semantic
normal form. An explicit simp list makes the dependency visible, and the
unchanged theorem statement catches any resulting mismatch. The patch is
small enough that introducing a new object-specification tactic would add
more machinery than it removes.

\section{Factor the Pell comparison through a positive scale}
\label{sec:pell}

\subsection{Keep the native denominator normalization visible}

The ordinary ratio proof in
\hsrc{Paper1980/PellRatio90}{\code{Jones1980.Ratio90}} and the native proof in
\src{NativeHalfBinomial}{\code{upper_estimate_scaled}} use the same Pell growth
argument. The native source has a first coordinate twice the Pell sequence
value. That factor must remain in the statement and in the final ratio.

Write, with $U,Y,r\geq1$,
\[
 J=2r+1,\qquad a=Y(U+1),\qquad A=a+2,\qquad P=2UY^2+1,
\]
\[
 c=\psi_A(J),\qquad k_0=\psi_P(r+1),\qquad
 \xi=\frac{(U+1)^{2r}}{U^r}.
\]
The common theorem, named \code{upper_estimate_budget}, accepts a positive
real scale $\lambda$ (spelled \code{scale} in Lean). Its distinctive hypothesis
and conclusion are
\[
 \xi<\lambda(Y+1)
 \quad\Longrightarrow\quad
 \frac{c}{k_0}<\xi+\frac{\lambda}{2},
\]
under the existing growth assumptions $72r<U+1$ and $16\leq Y(U+1)$, and the
same Pell identities and positivity premises as the old proof.

\subsection{Why the generalization works}

Set $\varepsilon=9/[Y(U+1)]$. The shared analytic work establishes
\[
 \frac{c}{k_0}\leq\xi(1+2r\varepsilon),
 \qquad (Y+1)\,2r\varepsilon<\frac12.
\]
The first statement follows from the existing upper/lower growth estimates,
the polynomial comparison \code{upper_key}, and
\code{one_add_pow_le}. The second follows from the same numerical guard
$72r<U+1$. Since $\lambda>0$ and $2r\varepsilon\geq0$,
\[
 \xi\,2r\varepsilon
 \leq\lambda(Y+1)\,2r\varepsilon
 <\frac\lambda2.
\]
The final upper bound is then a single algebraic conclusion. The long
power-and-denominator argument is proved once; the normalization scale enters
only in the last error estimate.

Specializing to $\lambda=1$ gives the old ordinary
\code{upper_estimate'} statement. Specializing to $\lambda=2$ gives the old
native \code{upper_estimate_scaled} statement. The existing doubled-coordinate
consumer then uses $k=2k_0$ to obtain
\[
 \frac{c}{k}<\frac\xi2+\frac12.
\]
All original public arguments remain accepted, including an ordinary-wrapper
premise that is not needed by the new proof. This avoids a gratuitous
downstream signature migration.

\subsection{Reuse the lower comparison at the same boundary}

The native \src{NativePellPower}{\code{lower_key}} has the exact statement of
the lower-level \code{Ratio90.lower_key}; its local polynomial proof can be an
application. More significantly, \code{doubled_ratio_lower} follows from
\code{Ratio90.lower_estimate} applied at $k_0$, multiplying by two and rewriting
the native equation $k=2k_0$. Its current public statement is preserved.

The new import goes from \code{NativePellPower} to the existing H10 ratio
module. That module has no return dependency on \code{DiophantineConsistency}.
The full overlay check rebuilds the intervening native scale modules before
checking \code{NativeHalfBinomial}, so validation does not accidentally use
an old compiled native theorem beneath a changed dependency.

\subsection{Cost and relation to prior work}

The three affected complete modules shrink by \RatioRawSaved{} lines against
the pin, counting the helper once and the new import. Earlier reports already
suggested replacing existential aliases for $\xi,X,Z,\varepsilon$ with
\code{set} or \code{let}. After charging the valid cleanup, including its
redundant rewrite removal, to the baseline, the additional saving is
\RatioSaved{} lines. Current source already includes the two-line cleanup
inside the native \code{lower_key}; that reduction is already in the baseline
and is not deducted twice. The source and counterfactual metrics are supplied
separately so that this small overlap is not hidden.

The contribution is the scale-parametric inequality and its two mathematical
specializations, together with exact lower-comparison reuse. It does not
replace the earlier Pell-register decoding abstraction. The main review
question is whether this common theorem belongs in the shared H10 module;
the dependency direction and both existing clients support that placement.
No sharper constants or broader optimality claim is made.

\section{Optional: isolate integer recovery from two error budgets}

\subsection{A shared discrete conclusion}

The end of \src{NativeHalfBinomial}{\code{half_binomial_of_ratio}} combines a
binomial decomposition, an evenness argument, and two real inequalities to
identify a natural number. The ordinary
\hsrc{Paper1980/PellRatio90}{\code{Ratio90.Y_eq_floor}} has the same final
discrete task with different error constants.

The checked example separates that task. Let $Y,t\in\mathbb N$ and
$\rho,\delta,\epsilon\in\mathbb R$, and assume
\[
 \delta\geq0,\qquad\delta+\epsilon\leq1,
\]
\[
 Y<\rho<Y+1,\qquad t+\delta<\rho<t+\delta+\epsilon.
\]
Then $Y=t$. The second certificate implies $t\leq\rho<t+1$, so
\code{Nat.floor} of $\rho$ is both $Y$ and $t$. The proof uses
\mathsrc{Mathlib/Algebra/Order/Floor/Semiring.lean}{\code{Nat.floor_eq_iff}}
twice and transports equality; it does not reprove properties of floors.

This interface exposes which analytic facts actually matter to integer
recovery. A separate assumption $\epsilon\geq0$ is unnecessary: the strict
interval premises already imply a positive width. The proof does need
nonnegativity of $\delta$, and its total-error bound must be charged to the
same unit interval.

\subsection{Checked applications and limits}

\path{examples/IntervalExtraction.lean} contains the core and three checked
applications:

\begin{enumerate}
\item Native half-binomial recovery takes $\delta=f/2$ and $\epsilon=1/2$
  with $0\leq f<1$. The zero endpoint is allowed, matching the actual
  decomposition theorem.
\item Ordinary binomial recovery takes $0\leq\delta<1/4$ and
  $\epsilon=1/2$.
\item A variable-denominator extension takes $s\geq2$,
  $\delta=f/s$, and $\epsilon=1/s$, again with $0\leq f<1$.
\end{enumerate}

The budget is substantive. If it is allowed to exceed one, $t=0$, $Y=1$,
$\rho=5/4$, $\delta=0$, and $\epsilon=3/2$ satisfy both strict intervals but
give unequal integers. If negative $\delta$ is allowed, $t=1$, $Y=0$,
$\rho=1/2$, $\delta=-3/4$, and $\epsilon=1/2$ give another counterexample.
These examples explain the contract; they are not new conditions added to
the existing clients, which already meet it.

The actual parity proof, divisibility by two, and real decomposition remain
the clients' responsibility. The supplied file checks the final
client-shaped implications; it does not replace the complete production
modules with a new rounding API. For only two consumers, the common lemma
and adapters may add lines. Its value is a stable mathematical boundary for
further normalization variants, so no net saving enters the total.

\section{Optional: share closure of observed internal certificate sets}

\subsection{The repeated operation is union and insertion}

\src{ObjectCodeBound}{\code{boundClosed_cup}} and
\code{boundClosed_insert} establish closure when two internal certificate sets
are united or one justified observation is adjoined. The imported
\fullsrc{SetTheory/BoundedConsistency/Lean/BoundedZFCConsistency/InternalSoundness.lean}{\code{InternalSoundness}}
renaming certificates and
\fullsrc{SetTheory/BoundedConsistency/Lean/BoundedZFCConsistency/CodedRank.lean}{\code{CodedRank}}
rank certificates repeat that pattern. Their local step predicates differ,
but each is monotone as the supporting internal set grows.

Use an observation type $I$, the arbitrary model carrier $V$, and its
membership relation. The interface stores an encoder $e:I\to V$, a step
predicate $\operatorname{step}(S,i)$, and monotonicity under internal set
inclusion. Define
\[
 \operatorname{Closed}(S)
 \;\Longleftrightarrow\;
 \forall i,\ e(i)\in S\ \Longrightarrow\ \operatorname{step}(S,i).
\]
For bound certificates, $I=V\times V$, $e(c,b)$ is the existing ordered-pair
encoding, and the step predicate includes both $b\in\omega$ and the existing
\code{BoundStep}. No guard is discarded to obtain a cleaner abstraction.

\subsection{The two laws and their necessary hypotheses}

If $S$ and $T$ are closed, an observed member of $S\cup T$ lies in one of
the inputs. The corresponding step follows from that input's closure and
then transports along its inclusion into the union. This proves union
closure once for every rule family.

For insertion, assume $e$ injective, $S$ closed, and
$\operatorname{step}(S,i)$. Consider an observation in
$S\cup\{e(i)\}$. An old member is handled by monotonicity. A member of the
new singleton has encoding equal to $e(i)$; injectivity identifies its
observation with $i$, and the supplied step transports to the enlarged set.

For unrestricted step predicates, injectivity is essential to this insertion
interface. If two different observations have the same
encoding, adjoining that encoding can create an obligation for the other
observation, whose step was never supplied. The concrete adapters prove
injectivity using the repository's \code{kpair_inj} or \code{satTriple_inj};
the generic theorem does not silently assume it.

The encoding also need not cover the carrier. Unrelated or malformed members
of an internal set are allowed, just as in the original closure definitions:
only encoded observations trigger obligations. Requiring every member to be
a well-formed pair would strengthen the certificate and change the public
meaning. Replacing the model's internal set with an external \code{Finset}
would likewise be unjustified.

\subsection{What has actually been checked}

\path{examples/ObservedClosure.lean} supplies the kernel, including the empty
set, union, and insertion. \path{examples/ClosureClients.lean} supplies six
concrete client theorems: union and insertion for bound, renaming, and rank
certificates. Both complete example files compile against the pinned
repository's actual foundational definitions.

Only the bound family is in the requested folder. The renaming and rank
families are explicit possible amortization sites in imported foundational
code; they are not silently counted as reductions in
\code{DiophantineConsistency}. The production patch does not introduce this
architecture, and no net source reduction is claimed. Adoption becomes
attractive when several families use the kernel and stop maintaining their
own insertion arguments.

The prior sequential scanner-fold abstraction has a changing accepted-history
state and list execution. The present laws concern unordered internal sets
and their union/singleton operations. Closure alone does not establish
well-foundedness, termination, or correctness of certificates; the existing
internal induction and constructor-inversion arguments keep those roles.

\section{Validation, reproduction, and adoption order}

\subsection{Evidence and its scope}

The compiler is the repository's Lean 4.32.0. The mathlib commit is
\path{81a5d257c8e410db227a6665ed08f64fea08e997}; the validation scripts check
every dependency revision against \code{lake-manifest.json}. Repository
prerequisites were built from the pinned sources in an isolated output tree.
The final common overlay then checked all changed production modules and
their affected local dependencies together.

The final overlay's import closure contains \CombinedModules{} repository
modules. Its record distinguishes newly compiled outputs from unchanged
baseline outputs reused only after source, dependency-object, and output
hashes matched. This is an import-closure check, not a rebuild of every target
in the repository. The package dependencies use their already-built artifacts
at the verified revisions.

For the specification family, 16 additional full-source audit compilations
printed the axioms of the 20 changed theorems before and after. Each retained
exactly \code{propext}, \code{Classical.choice}, and \code{Quot.sound}. The
other new examples and the final integrated endpoints were also checked with
\code{#print axioms}; their reported dependencies contain no \code{sorryAx}
or newly postulated project axiom. Ordinary Lean foundational axioms are not
described as an axiom-free proof.

The source metrics count physical lines, including imports, comments, and
blank lines. They are reproducible for the shipped candidates. The
counterfactual Pell adjustment is narrower: it applies the known prior alias
cleanup before calculating the additional reduction. Earlier architectural
packets contain alternatives and prospective designs, so this report does
not invent a globally combined post-proposal checkout. The novelty claim is
about the new mechanisms and obligations, with that limitation made explicit.

Elapsed times in raw logs are diagnostic data only. Shared caches, concurrent
checks, and local resource conditions make them unsuitable for a performance
comparison. An interrupted baseline build during disk-space recovery was
rerun from regenerated outputs; the final successful records identify the
artifacts used in the overlay.

\subsection{A concrete integration sequence}

\begin{enumerate}
\item Apply the formula-specification patch. It changes proof bodies only,
  has no helper cost, and gives the largest immediate reduction.
\item Review the common Pell theorem together with all four public adapters.
  Keep the factor of two explicit and rebuild native scale consumers.
\item Revisit the two optional abstractions when there are enough clients to
  amortize their contracts, encoders, and adapters. Do not count two-line
  applications while ignoring the new abstraction's source.
\end{enumerate}

The patches target the pinned checkout. If earlier proposals have already
changed the same lines in a working branch, port the mechanism into their
resulting interfaces and rerun the supplied checks; do not overwrite those
changes with the full candidate snapshots. The snapshots exist to make this
review exact and independently reproducible.

\subsection{Commands and artifact map}

The included \path{scripts/validate.py} constructs an isolated source overlay
without changing the repository. It checks the repository pin, Lean version,
and dependency commits; discovers local imports; and compiles the selected
closure. Dependency package artifacts must already be available, as for a
normal local Lean project build. A typical invocation is:

\begin{lstlisting}[language={}]
python3 scripts/validate.py \
  --repo /path/to/ProveIt \
  --lean /path/to/lean-4.32.0/bin/lean \
  --packages /path/to/ProveIt/.lake/packages \
  --out /tmp/diophantine-simplification-check

python3 scripts/check_examples.py \
  --repo /path/to/ProveIt \
  --lean /path/to/lean-4.32.0/bin/lean \
  --packages /path/to/ProveIt/.lake/packages \
  --validated /tmp/diophantine-simplification-check
\end{lstlisting}

The README gives the exact prerequisites and explains the optional verified
baseline reuse. The shipped evidence contains the actual successful run,
rather than claiming these portable commands built a full fresh dependency
cache. PDF reproduction uses two \code{pdflatex} passes on \path{report.tex}.

\begin{longtable}{@{}P{49mm}P{100mm}@{}}
\toprule
Artifact & Purpose\\
\midrule\endhead
\path{report.tex}, \path{report.pdf} & Editable manuscript and rendered report, with exact source links.\\
\path{metrics.tex} & Generated numerical values used by the manuscript.\\
\path{patches/combined.patch} & One reviewable patch for the two integrated directions.\\
\path{patches/} & Independent family patches for selective adoption.\\
\path{candidates/} & Complete proposed production sources, preserving repository-relative paths.\\
\path{examples/} & Checked optional cores and small independent demonstrations.\\
\path{evidence/} & Novelty audit, inventories, source metrics, compiler records, and axiom reports.\\
\path{scripts/} & Pin-aware compilation, example/axiom checks, and integrity/metric verification.\\
\path{validation-plan.json} & Exact production endpoints and example/audit targets.\\
\path{SHA256SUMS} & Integrity hashes for all other files in the packet.\\
\bottomrule
\end{longtable}

\appendix
\section{Exact specification targets}

The proof-only patch changes the following names in their existing module
namespaces. Original source ranges and before/after hashes are recorded in
the machine-readable theorem and module metrics.

\small
\begin{longtable}{@{}P{55mm}P{94mm}@{}}
\toprule
Module & Target declarations\\
\midrule\endhead
\code{ObjectAdditionLaws} & \code{successorLeftStage_spec}, \code{commutativeStage_spec}, \code{associativeStage_spec}\\
\code{ObjectAdditionCancellation} & \code{cancellationStage_spec}\\
\code{ObjectMultiplicationLaws} & \code{successorLeftStage_spec}, \code{commutativeStage_spec}\\
\code{ObjectDistributivity} & \code{distributiveStage_spec}, \code{associativeStage_spec}\\
\code{ObjectNaturalOrder} & \code{additionBoundStage_spec}, \code{differenceStage_spec}\\
\code{ObjectNaturalExpression} & \code{evalForm_spec}\\
\code{ObjectCodeBound} & \code{boundAtomForm_spec}, \code{boundBinaryForm_spec}, \code{boundQuantifierForm_spec}, \code{boundStepForm_spec}, \code{boundClosedForm_spec}, \code{boundCodeForm_spec}\\
\code{ObjectCodeSeal} & \code{sealCodeForm_spec}, \code{sealedZFCAxiomCodeForm_spec}, \code{ctxSealedZFCAxiomsForm_spec}\\
\bottomrule
\end{longtable}
\normalsize

\end{document}
